\newpage
\textbf{\huge CHAPTER  5}\\
\section{ONGOING RESULTS AND COVERAGE IMPACT} \setcounter{section}{5}
\subsection{Introduction to Performance Metrics}\setcounter{subsection}{1}
% \pagenumbering{arabic} 
\begin{justify}
    As the migration from Apache Storm to Apache Flink is a highly critical infrastructure upgrade for Flipkart's Search team, evaluating the success of the project relies on rigorous performance benchmarking. The primary metrics for evaluation are ingestion latency (specifically the p99 latency), system throughput under heavy load, resource utilization (CPU and memory footprint), and strict data parity between the legacy and modern pipelines. 

    At the current stage of the internship, the Flink pipeline has been successfully developed and deployed in a shadow environment. This chapter discusses the ongoing results obtained from this shadow deployment and the projected coverage impact on the live Search ecosystem.
\end{justify}

\vspace{0.25cm}

\subsection{Shadow Testing and Data Parity}
\begin{justify}
    Deploying a new real-time ingestion engine directly to production poses severe risks to Flipkart's business operations. Therefore, a "Shadow Testing" approach was utilized. The Flink topology was configured to consume from the same production Kafka source topics as the live Storm topology. However, Flink's output was routed to isolated sink topics and a staging Solr cluster.
    
    To guarantee exactly-once processing and correct business logic translation, comprehensive automated data validation scripts were executed. These scripts continuously compared the final product entities generated by Flink against those generated by Storm. Over a testing period encompassing millions of catalog and price updates, the Flink pipeline achieved \textbf{100\% data parity}. This successfully validated that the complex state merging logic and Base View Transformations were accurately ported to the new framework without any loss of business context.
\end{justify}

\subsection{Latency and Throughput Improvements}
\begin{justify}
    The most profound result observed during the shadow testing phase was the drastic reduction in processing latency, primarily driven by Flink's localized state management (RocksDB). 
    
    In the legacy Storm architecture, every state update required a synchronous network call to the Apache HBase cluster. During simulated peak load testing (mimicking Big Billion Days traffic), Storm's external network I/O became a significant bottleneck, causing Kafka consumer lag and increasing the p99 processing latency.
    
    Conversely, the Flink pipeline, utilizing its embedded RocksDB state backend, eliminated these external lookups. State retrieval and updates were executed locally on the TaskManager's disk in sub-milliseconds. 
    
    \begin{itemize}
        \item \textbf{Throughput:} Flink demonstrated a significantly higher throughput capacity, easily processing tens of thousands of messages per second per node without triggering backpressure.
        \item \textbf{Latency:} The end-to-end ingestion latency (from Kafka source to Kafka sink) saw a projected reduction of over 40\% at the p99 percentile, ensuring that critical signals like "Out of Stock" reach the Search Index significantly faster.
    \end{itemize}
\end{justify}

\subsection{Resource Utilization and Infrastructure Efficiency}
\begin{justify}
    Beyond speed, the migration to Apache Flink showcased substantial improvements in infrastructure efficiency. Apache Storm's architecture requires over-provisioning of worker nodes to handle peak traffic spikes and the overhead of constant network I/O. 

    By leveraging Flink's optimized execution graph and localized state, the new pipeline achieved the same, if not better, throughput utilizing fewer computational resources. The CPU utilization across the Flink cluster was observed to be more stable, and memory management was handled efficiently by Flink's native memory manager off-heap. This translates directly to reduced cloud infrastructure costs for the Search team while providing a higher ceiling for future scalability.
\end{justify}

\subsection{Projected Coverage Impact on Search Ecosystem}
\begin{justify}
    The engineering improvements brought by Flink have a direct, measurable impact on Flipkart's business metrics and overall customer experience. The "Coverage Impact" refers to how accurately and quickly the Search Index reflects the real-world state of the catalog.
    
    \begin{enumerate}
        \item \textbf{Reduced Price Anomaly:} Faster ingestion of price drop signals ensures that customers searching for products see the most updated, competitive prices instantaneously, improving click-through rates (CTR).
        \item \textbf{Inventory Accuracy:} NRT processing of availability signals minimizes the window where an out-of-stock product remains visible in search results. This directly reduces order cancellations and customer frustration.
        \item \textbf{System Resilience:} Flink's automated checkpointing and exactly-once semantics ensure that even during partial cluster failures, no update signals are dropped. This guarantees that the Search Index maintains 100\% eventual consistency with the upstream business databases without manual intervention.
    \end{enumerate}
\end{justify}

\begin{table}[h]
    \centering
    \small
    \caption{Performance Comparison: Storm vs. Flink (Shadow Environment)}
    \renewcommand{\arraystretch}{1.5}
    \begin{tabular}{|p{4cm}|p{5cm}|p{5cm}|}
        \hline
        \textbf{Metric} & \textbf{Legacy System (Apache Storm)} & \textbf{Proposed System (Apache Flink)} \\
        \hline
        State Management & External (HBase) - High Latency & Localized (RocksDB) - Sub-ms Latency \\
        \hline
        Processing Semantics & At-least-once & Exactly-once \\
        \hline
        Peak Load Behavior & High consumer lag, Backpressure & Stable throughput, Low lag \\
        \hline
        Data Parity Status & Baseline & 100\% Matched (Validated) \\
        \hline
    \end{tabular}
    \label{tab:performance_metrics}
\end{table}